On considère l’équation différentielle (E) : $2y' + y = e^{-0,5x}$ où $y$ est une fonction de la variable réelle $x$, définie et dérivable sur l’ensemble $\R$ des nombres réels, et $y'$ est la fonction dérivée de $y$.

\begin{enumerate}
	\item Donner et résoudre l'équation homogène (E$_0$)\sautp

	\Lignes{2}
	\item Démontrer, par un calcul, que la fonction $p$ définie sur $\R$ par : $p(x) = 0,5~x~e^{-0,5x}$ est une solution particulière de l'équation différentielle (E)\sautp

	\Lignes{6}
	\item En déduire la forme générale des solutions de (E).\sautp

	\Lignes{1}
	\item Déterminer, à l'aide de géogebra et d'un curseur bien choisi, la solution particulière $f$ de (E) telle que $f(0)=1$.\sautp

	\Lignes{1}
	
	\appelprof
\end{enumerate}